\chapter{Mini-Examenes de Trigonometria}

\section{[C05.S01] Grados y radianes}

\begin{exambrief}{Mini-examen C05.S01}
Duracion sugerida: \(25\) minutos. Valor orientativo: \(10\) puntos.
\begin{enumerate}
  \item Convierte \(-\frac{3\pi}{4}\) radianes a grados y \(315^\circ\) a radianes.
  \item Una antena parte orientada a \(40^\circ\), gira \(\frac{7\pi}{6}\) radianes en sentido
  antihorario y corrige despues \(95^\circ\) en sentido horario. Halla la orientacion final en
  grados y en radianes.
\end{enumerate}
\end{exambrief}

\begin{shortanswer}{Mini-examen C05.S01}
\[
-\frac{3\pi}{4}=-135^\circ,
\qquad
315^\circ=\frac{7\pi}{4}.
\]
La orientacion final es
\[
155^\circ=\frac{31\pi}{36}.
\]
\end{shortanswer}

\begin{fullsolution}{Mini-examen C05.S01}
\[
-\frac{3\pi}{4}\cdot \frac{180^\circ}{\pi}=-135^\circ,
\qquad
315^\circ\cdot \frac{\pi}{180^\circ}=\frac{7\pi}{4}.
\]
Para la antena,
\[
\frac{7\pi}{6}\text{ rad}=210^\circ.
\]
Luego el giro total es
\[
40^\circ+210^\circ-95^\circ=155^\circ.
\]
En radianes,
\[
155^\circ=\frac{155\pi}{180}=\frac{31\pi}{36}.
\]
\end{fullsolution}

\section{[C05.S02] Razones trigonometricas en triangulos y poligonos}

\begin{exambrief}{Mini-examen C05.S02}
Duracion sugerida: \(30\) minutos. Valor orientativo: \(10\) puntos.
\begin{enumerate}
  \item Una escalera de \(\SI{7}{m}\) forma \(65^\circ\) con el suelo. Halla la altura alcanzada
  y la distancia horizontal a la pared.
  \item Un hexagono regular tiene lado \(\SI{6}{m}\). Calcula su apotema y su area.
\end{enumerate}
\end{exambrief}

\begin{shortanswer}{Mini-examen C05.S02}
\[
h=7\sen 65^\circ\approx \SI{6.34}{m},
\qquad
d=7\cos 65^\circ\approx \SI{2.96}{m}.
\]
Para el hexagono,
\[
a=3\sqrt3,
\qquad
A=54\sqrt3\ \text{m}^2.
\]
\end{shortanswer}

\begin{fullsolution}{Mini-examen C05.S02}
En el triangulo rectangulo de la escalera:
\[
h=7\sen 65^\circ\approx 7\cdot 0.9063\approx \SI{6.34}{m},
\]
\[
d=7\cos 65^\circ\approx 7\cdot 0.4226\approx \SI{2.96}{m}.
\]
En el hexagono regular, el apotema es la altura de un triangulo equilatero de lado \(6\):
\[
a=\frac{6\sqrt3}{2}=3\sqrt3.
\]
El area de un triangulo equilatero de lado \(6\) es
\[
\frac{6^2\sqrt3}{4}=9\sqrt3.
\]
Como hay seis triangulos:
\[
A=6\cdot 9\sqrt3=54\sqrt3.
\]
\end{fullsolution}

\section{[C05.S03] Angulos notables y reduccion al primer cuadrante}

\begin{exambrief}{Mini-examen C05.S03}
Duracion sugerida: \(25\) minutos. Valor orientativo: \(10\) puntos.
\begin{enumerate}
  \item Calcula exactamente \(\cos 150^\circ\) y \(\sen 330^\circ\).
  \item Un cable de \(\SI{8}{m}\) apunta con direccion \(225^\circ\). Halla exactamente las
  coordenadas de su extremo.
\end{enumerate}
\end{exambrief}

\begin{shortanswer}{Mini-examen C05.S03}
\[
\cos 150^\circ=-\frac{\sqrt3}{2},
\qquad
\sen 330^\circ=-\frac12.
\]
El extremo del cable es
\[
(-4\sqrt2,-4\sqrt2).
\]
\end{shortanswer}

\begin{fullsolution}{Mini-examen C05.S03}
\[
\cos 150^\circ=-\cos 30^\circ=-\frac{\sqrt3}{2},
\qquad
\sen 330^\circ=-\sen 30^\circ=-\frac12.
\]
Como \(225^\circ=180^\circ+45^\circ\),
\[
\cos 225^\circ=-\frac{\sqrt2}{2},
\qquad
\sen 225^\circ=-\frac{\sqrt2}{2}.
\]
Las coordenadas del extremo del cable son
\[
(8\cos225^\circ,8\sen225^\circ)=(-4\sqrt2,-4\sqrt2).
\]
\end{fullsolution}

\section{[C05.S04] Calculo de razones a partir de una razon conocida}

\begin{exambrief}{Mini-examen C05.S04}
Duracion sugerida: \(25\) minutos. Valor orientativo: \(10\) puntos.
\begin{enumerate}
  \item Sabiendo que \(\tg\alpha=-\frac34\) y que \(\alpha\) esta en el segundo cuadrante,
  calcula \(\sen\alpha\) y \(\cos\alpha\).
  \item Si un haz de longitud \(\SI{15}{m}\) forma ese mismo angulo \(\alpha\), halla las
  coordenadas de su extremo.
\end{enumerate}
\end{exambrief}

\begin{shortanswer}{Mini-examen C05.S04}
\[
\sen\alpha=\frac35,
\qquad
\cos\alpha=-\frac45.
\]
El extremo del haz es
\[
(-12,9).
\]
\end{shortanswer}

\begin{fullsolution}{Mini-examen C05.S04}
Si \(\tg\alpha=-\frac34\) en el segundo cuadrante, el seno es positivo y el coseno negativo.
Tomamos un triangulo de referencia con catetos \(3\) y \(4\), de hipotenusa \(5\). Entonces
\[
\sen\alpha=\frac35,
\qquad
\cos\alpha=-\frac45.
\]
Las coordenadas del extremo del haz de longitud \(15\) son
\[
(15\cos\alpha,15\sen\alpha)=\left(15\cdot -\frac45,15\cdot \frac35\right)=(-12,9).
\]
\end{fullsolution}

\section{[C05.S05] Identidades y transformaciones trigonometricas}

\begin{exambrief}{Mini-examen C05.S05}
Duracion sugerida: \(30\) minutos. Valor orientativo: \(10\) puntos.
\begin{enumerate}
  \item Simplifica
  \[
  E(x)=\frac{\sen 2x}{1+\cos 2x}.
  \]
  \item Resuelve \(E(x)=\sqrt3\) en \([0,2\pi)\).
\end{enumerate}
\end{exambrief}

\begin{shortanswer}{Mini-examen C05.S05}
\[
E(x)=\tg x.
\]
Las soluciones son
\[
x=\frac{\pi}{3},\qquad x=\frac{4\pi}{3}.
\]
\end{shortanswer}

\begin{fullsolution}{Mini-examen C05.S05}
Usamos
\[
\sen 2x=2\sen x\cos x,
\qquad
1+\cos 2x=2\cos^2 x.
\]
Entonces
\[
E(x)=\frac{2\sen x\cos x}{2\cos^2x}=\tg x.
\]
Resolvemos
\[
\tg x=\sqrt3.
\]
El angulo de referencia es \(\frac{\pi}{3}\) y la tangente es positiva en los cuadrantes I y
III. Por tanto,
\[
x=\frac{\pi}{3},\qquad x=\pi+\frac{\pi}{3}=\frac{4\pi}{3}.
\]
\end{fullsolution}

\section{[C05.S06] Ecuaciones trigonometricas}

\begin{exambrief}{Mini-examen C05.S06}
Duracion sugerida: \(30\) minutos. Valor orientativo: \(10\) puntos.
\begin{enumerate}
  \item Resuelve en \([0,2\pi)\)
  \[
  2\sen^2 x-3\sen x+1=0.
  \]
  \item Resuelve en \([0,2\pi)\)
  \[
  2\cos x=-\sqrt3.
  \]
\end{enumerate}
\end{exambrief}

\begin{shortanswer}{Mini-examen C05.S06}
\[
x=\frac{\pi}{6},\ \frac{\pi}{2},\ \frac{5\pi}{6}
\]
para la primera ecuacion, y
\[
x=\frac{5\pi}{6},\ \frac{7\pi}{6}
\]
para la segunda.
\end{shortanswer}

\begin{fullsolution}{Mini-examen C05.S06}
En la primera ecuacion hacemos \(u=\sen x\):
\[
2u^2-3u+1=0=(2u-1)(u-1).
\]
Por tanto,
\[
\sen x=\frac12 \quad \text{o}\quad \sen x=1.
\]
En \([0,2\pi)\),
\[
\sen x=\frac12 \Longrightarrow x=\frac{\pi}{6},\ \frac{5\pi}{6},
\qquad
\sen x=1 \Longrightarrow x=\frac{\pi}{2}.
\]

En la segunda:
\[
\cos x=-\frac{\sqrt3}{2}.
\]
Eso ocurre en los cuadrantes II y III:
\[
x=\frac{5\pi}{6},\ \frac{7\pi}{6}.
\]
\end{fullsolution}

\section{[C05.S07] Resolucion de triangulos}

\begin{exambrief}{Mini-examen C05.S07}
Duracion sugerida: \(30\) minutos. Valor orientativo: \(10\) puntos.
\begin{enumerate}
  \item En un triangulo se conocen dos lados de \(\SI{10}{m}\) y \(\SI{13}{m}\), y el angulo
  comprendido de \(47^\circ\). Halla el tercer lado.
  \item Calcula el area del triangulo.
\end{enumerate}
\end{exambrief}

\begin{shortanswer}{Mini-examen C05.S07}
El tercer lado mide aproximadamente \(\SI{9.58}{m}\) y el area es aproximadamente
\(\SI{47.6}{m^2}\).
\end{shortanswer}

\begin{fullsolution}{Mini-examen C05.S07}
Por el teorema del coseno,
\[
c^2=10^2+13^2-2\cdot 10\cdot 13\cos 47^\circ.
\]
Entonces
\[
c^2=269-260\cos 47^\circ\approx 269-177.32=91.68.
\]
Por tanto,
\[
c\approx \sqrt{91.68}\approx \SI{9.58}{m}.
\]
El area vale
\[
A=\frac12 ab\sen C=\frac12\cdot 10\cdot 13\sen 47^\circ
=65\sen 47^\circ\approx \SI{47.6}{m^2}.
\]
\end{fullsolution}

\section{[C05.S08] Problemas contextualizados de trigonometria}

\begin{exambrief}{Mini-examen C05.S08}
Duracion sugerida: \(35\) minutos. Valor orientativo: \(10\) puntos.
\begin{enumerate}
  \item Dos puntos de observacion estan sobre una carretera recta y separados \(\SI{30}{m}\).
  Desde el mas alejado de la base de una antena se ve su cima con angulo \(35^\circ\), y desde
  el mas cercano con angulo \(50^\circ\).
  \item Calcula la distancia desde el punto mas cercano hasta la base y la altura de la antena.
\end{enumerate}
\end{exambrief}

\begin{shortanswer}{Mini-examen C05.S08}
La distancia desde el punto cercano hasta la base es aproximadamente \(\SI{42.7}{m}\) y la
altura de la antena es aproximadamente \(\SI{50.9}{m}\).
\end{shortanswer}

\begin{fullsolution}{Mini-examen C05.S08}
Sea \(x\) la distancia desde el punto cercano hasta la base. Entonces
\[
\tg 50^\circ=\frac{h}{x},
\qquad
\tg 35^\circ=\frac{h}{x+30}.
\]
Igualando:
\[
x\tg50^\circ=(x+30)\tg35^\circ.
\]
Despejamos:
\[
x(\tg50^\circ-\tg35^\circ)=30\tg35^\circ,
\]
\[
x=\frac{30\tg35^\circ}{\tg50^\circ-\tg35^\circ}
\approx \SI{42.7}{m}.
\]
La altura es
\[
h=x\tg50^\circ\approx 42.7\cdot 1.1918\approx \SI{50.9}{m}.
\]
\end{fullsolution}
